\section{From MAIA to MAIA-AV: Dataset Expansion and Curation}
\subsection*{Video Selection}
The dataset adopted in MAIA-AV builds upon the guidelines of \ac{MAIA}, comprising 100 short videos, approximately 30 seconds each, selected from YouTube Italia to represent everyday scenarios of Italian culture. Content domains include sports, culinary arts, urban life, and cultural activities, with a strong focus on human presence and highly informative scenes. In the original formulation, video selection and annotation were oriented strictly toward visual comprehension, completely omitting audio content.
Building upon this initial set, we collected an additional 100 videos, each uniquely identified from \texttt{video001} to \texttt{video100}. These identifiers map each video to its corresponding audio stream and text transcriptions. Each video was precisely trimmed to extract targeted segments, creating a controlled corpus of clips. Particular attention was paid to the alignment between dialogue content and the elements appearing in the scene.
This targeted selection of relevant dialogue enables a precise investigation of multimodal interactions. By prioritizing scenarios where spoken dialogue is semantically grounded in the visual content, we ensure that the audio signal plays a genuinely informative role in the multimodal instance.
This distinction represents the core departure from the original dataset: \ac{MAIA}-AV is a benchmark engineered to combine and compare distributed evidence across visual, auditory, and textual modalities.

\paragraph{Implementation Details}
The video acquisition pipeline relies on \texttt{yt\_dlp} to retrieve targeted temporal segments from selected YouTube videos. Rather than downloading complete sources for post-hoc editing, each clip is specified prior to acquisition via three parameters: source URL, start timestamp, and duration. These parameters configure a precise temporal interval directly governing the retrieval process.
Timestamp conversion into seconds is handled by \texttt{parse\_time}, which reduces \texttt{HH:MM:SS} or \texttt{MM:SS} strings into integer values:
\begin{lstlisting}[language=Python, xleftmargin=2em]
def parse_time(time_str):
    """Converts HH:MM:SS or MM:SS string to seconds."""
    parts = [int(p) for p in time_str.split(':')]
    
    if len(parts) == 3:
        return parts[0] * 3600 + parts[1] * 60 + parts[2]
    elif len(parts) == 2:
        return parts[0] * 60 + parts[1]
    return parts[0]
\end{lstlisting}
Adding the requested duration to the start position determines the upper boundary of the segment:
\begin{lstlisting}[language=Python, xleftmargin=2em]
start_sec = parse_time(start_time_str)
end_sec = start_sec + int(durata_str)
\end{lstlisting}
The resulting bounds, \texttt{start\_sec} and \texttt{end\_sec}, define the exact temporal window extracted from the source.
Output paths are structured prior to downloader configuration, ensuring MP4 container compliance and consistent dataset indexing under the target directory:
\begin{lstlisting}[language=Python, xleftmargin=2em]
filename_input = input(
    "Enter output filename (e.g., clip): "
).strip()
if not filename_input.endswith(".mp4"):
    filename_input += ".mp4"
output = f"Dataset/video{filename_input}"
\end{lstlisting}
Prefixing filenames with \texttt{video} maintains a unified identifier convention across all dataset instances.
Stream selection, temporal trimming, and media normalization are encapsulated within the \texttt{ydl\_opts} dictionary. Format selection prioritizes high-resolution video streams alongside optimal audio feeds:
\begin{lstlisting}[language=Python, xleftmargin=2em]
ydl_opts = {
    'format': 'bestvideo[height<=2160]+bestaudio/best',
}
\end{lstlisting}
This filter selects video streams up to 2160p paired with the highest quality audio stream, falling back to pre-merged streams (\texttt{/best}) if separate tracks are unavailable.
Temporal constraints are applied during acquisition via \texttt{download\_range\_func}:
\begin{lstlisting}[language=Python, xleftmargin=2em]
    'download_ranges': download_range_func(
        None,
        [(start_sec, end_sec)]
    ),
\end{lstlisting}
Embedding segment bounds directly into the retrieval mechanism eliminates separate media cropping steps and reduces bandwidth overhead.
To maintain frame-accurate cutting across keyframe-based compressed video formats, the pipeline enforces keyframe generation at boundary points:
\begin{lstlisting}[language=Python, xleftmargin=2em]
    'force_keyframes_at_cuts': True,
\end{lstlisting}
Forcing keyframes ensures temporal alignment at cut points, preserving action continuity for downstream multimodal evaluation in MAIA-AV.
Output formats and post-processing codecs are explicitly declared to standardize the corpus:
\begin{lstlisting}[language=Python, xleftmargin=2em]
    'merge_output_format': 'mp4',
    'postprocessor_args': {
        'ffmpeg': [
            '-c:v', 'libx264',
            '-preset', 'fast',
            '-c:a', 'aac'
        ]
    },
\end{lstlisting}
FFmpeg re-encodes visual streams via H.264 (\texttt{libx264}, \texttt{fast} preset) and acoustic streams via AAC, eliminating encoding heterogeneity across raw YouTube sources.
The unified configuration dictionary combines acquisition parameters, formatting rules, output paths, and authentication credentials:
\begin{lstlisting}[language=Python, xleftmargin=2em]
ydl_opts = {
    'format': 'bestvideo[height<=2160]+bestaudio/best',
    'download_ranges': download_range_func(
        None,
        [(start_sec, end_sec)]
    ),
    'force_keyframes_at_cuts': True,
    'merge_output_format': 'mp4',
    'postprocessor_args': {
        'ffmpeg': [
            '-c:v', 'libx264',
            '-preset', 'fast',
            '-c:a', 'aac'
        ]
    },
    'outtmpl': output,
    'cookiefile': 'src/youtube.com_cookies.txt',
}
\end{lstlisting}
Passing a session cookie file (\texttt{cookiefile}) enables retrieval of age-restricted or session-gated content without altering core extraction logic.
Execution is wrapped in exception handling blocks, logging errors and returning non-zero exit codes upon failure:
\begin{lstlisting}[language=Python, xleftmargin=2em]
try:
    with yt_dlp.YoutubeDL(ydl_opts) as ydl:
        ydl.download([url])
    print(f"\nSuccessfully downloaded clip to {output}")
except Exception as e:
    print(f"\nError: {str(e)}")
    sys.exit(1)
\end{lstlisting}
This setup allows \texttt{yt\_dlp} to coordinate format resolution, stream extraction, interval clipping, and FFmpeg transcoding within a single execution pass.

\subsection*{Audio Extraction}
In the MAIA-AV framework, the auditory element of each video is regarded as an autonomous source of information rather than a mere supplementary part of the visual narrative. Spoken language, ambient sounds, and acoustic phenomena offer significant evidence that enhances visual data and aids in scene interpretation. Isolating the audio allows this modality to be independently analyzed, ensuring its contribution is maintained throughout subsequent processing stages.
For every video, the original audio track is extracted and archived as a separate file. This separation retains the identical temporal parameters and identifier associated with the corresponding audiovisual entity. Such a one-to-one mapping upholds alignment across modalities, enabling visual, acoustic, and textual representations to be traced back to the same foundational event.
The resulting audio files function as direct inputs for further pipeline stages, including automatic speech transcription and specialized audio analyses. The extraction process neither generates a modified version of the original content nor alters it; instead, it separates the auditory channel, allowing independent processing while remaining completely synchronized with the originating video content.
\paragraph{Implementation Details}
The audio extraction stage converts each audiovisual dataset instance into a standalone acoustic representation while preserving multimodal alignment with the source video. Implemented via Python's \texttt{subprocess} module, the pipeline invokes FFmpeg to automate batch processing across the complete video collection.
Input and output paths are declared explicitly at execution setup:
\begin{lstlisting}[language=Python, xleftmargin=2em]
INPUT_FOLDER = Path("data/video")
OUTPUT_FOLDER = Path("data/audio")
VIDEO_EXTENSIONS = {".mp4"}
\end{lstlisting}
Decoupling input and output directories isolates raw video resources and establishes an aligned audio mirror directory. Restricting \texttt{VIDEO\_EXTENSIONS} to MP4 matches the output format produced by the preceding acquisition stage.
The core conversion routine is encapsulated within \texttt{mp4\_to\_mp3}:
\begin{lstlisting}[language=Python, xleftmargin=2em]
def mp4_to_mp3(
    input_path: Union[str, Path],
    output_path: Optional[Union[str, Path]] = None,
    bitrate: str = "192k",
) -> Path:
    """Extract the audio from an MP4 file and save it as MP3."""
\end{lstlisting}
When an explicit destination path is omitted, the function retains the base filename and replaces the container extension:
\begin{lstlisting}[language=Python, xleftmargin=2em]
    input_path = Path(input_path)
    output_path = (
        input_path.with_suffix(".mp3")
        if output_path is None
        else Path(output_path)
    )
\end{lstlisting}
This mapping strategy guarantees cross-modal alignment across representations (e.g., \texttt{video001.mp4} yields \texttt{video001.mp3}).
Prior to invoking FFmpeg, parent directory existence is enforced:
\begin{lstlisting}[language=Python, xleftmargin=2em]
    output_path.parent.mkdir(
        parents=True,
        exist_ok=True,
    )
\end{lstlisting}
This step allows nested input folder hierarchies to be automatically replicated within the output directory.
Extraction is executed using a single \texttt{subprocess.run} call:
\begin{lstlisting}[language=Python, xleftmargin=2em]
    subprocess.run(
        [
            "ffmpeg",
            "-y",
            "-i",
            str(input_path),
            "-vn",
            "-c:a",
            "libmp3lame",
            "-b:a",
            bitrate,
            str(output_path),
        ],
        check=True,
    )
\end{lstlisting}
The \texttt{-vn} flag discards video streams, ensuring that FFmpeg isolates the acoustic signal rather than performing generic media transcode operations.
Acoustic streams are normalized using the \texttt{libmp3lame} codec at a constant target bitrate of 192 kbit/s (\texttt{-b:a 192k}). Passing \texttt{-y} overwrites pre-existing destination files without prompt interruption, facilitating seamless pipeline re-execution. Setting \texttt{check=True} ensures runtime errors raise a \texttt{CalledProcessError} rather than failing silently.
Upon completion, the function returns the confirmed target path:
\begin{lstlisting}[language=Python, xleftmargin=2em]
    return output_path
\end{lstlisting}
Batch processing begins by recursively indexing and sorting compatible media files:
\begin{lstlisting}[language=Python, xleftmargin=2em]
def find_videos(folder: Path) -> list[Path]:
    """Return all supported video files found recursively."""
    return sorted(
        path
        for path in folder.rglob("*")
        if path.is_file()
        and path.suffix.lower() in VIDEO_EXTENSIONS
    )
\end{lstlisting}
Sorting file paths guarantees deterministic execution order across operating environments.
Directory validity is verified before initializing conversion loops:
\begin{lstlisting}[language=Python, xleftmargin=2em]
if not INPUT_FOLDER.exists():
    raise FileNotFoundError(
        f"Input folder not found: {INPUT_FOLDER}"
    )
videos = find_videos(INPUT_FOLDER)
if not videos:
    print(f"No videos found in: {INPUT_FOLDER}")
    return
\end{lstlisting}
Validating the input directory early avoids redundant FFmpeg invocations when processing empty or invalid paths.
For each video, relative path resolution preserves source subdirectory structures under \texttt{OUTPUT\_FOLDER}:
\begin{lstlisting}[language=Python, xleftmargin=2em]
for index, video_path in enumerate(videos, start=1):
    relative_path = video_path.relative_to(INPUT_FOLDER)
    audio_path = (
        OUTPUT_FOLDER
        / relative_path.with_suffix(".mp3")
    )
\end{lstlisting}
Replicating source directory trees ensures structural parity between visual and acoustic datasets.
Each file conversion is executed inside localized exception blocks:
\begin{lstlisting}[language=Python, xleftmargin=2em]
    try:
        mp4_to_mp3(
            video_path,
            audio_path,
        )
        print("Audio extraction completed.")
    except subprocess.CalledProcessError as error:
        print(f"FFmpeg error: {error}")
\end{lstlisting}
Capturing \texttt{CalledProcessError} exceptions locally prevents single-file execution failures from interrupting broader batch operations while logging specific processing errors.
\subsection*{Audio-Transcription Pipeline}
Another key representation in our dataset is the audio transcription, following the methodology established in our previous work \cite{deodati2026lost}. Transcriptions were generated using the \textit{Faster-Whisper} model (\texttt{large-v3}), a reimplementation of \textit{OpenAI's Whisper} powered by the \textit{CTranslate2} fast inference engine for Transformer architectures. The model processes raw audio content directly, with the language parameter set to Italian. Additionally, a \textit{Voice Activity Detection} (VAD) filter is applied to isolate segments containing actual speech, thereby minimizing non-speech transcriptions.
The overall transcription procedure aims to convert spoken dialogue into textual content that can be reliably recognized and processed by downstream models. This methodology represents an evolution compared to previous approaches on acoustic content analysis \cite{deodati2026lost}, which jointly employed \textit{Whisper-1} and \textit{GPT-4o-transcribe} to compare and align outputs for higher precision. In contrast, the current pipeline streamlines the transcription workflow by deploying a single, locally executed model, providing an efficient and reproducible framework.
\paragraph{Implementation Details}
The transcription stage was implemented as an independent component of the data preparation pipeline to associate each video with a textual representation of its spoken content. Powered by \texttt{Faster-Whisper} and the \texttt{large-v3} model executed locally on GPU, this setup ensures full reproducibility while eliminating dependencies on external cloud services.
The execution environment explicitly configures a selected CUDA device and establishes a local cache for model files:
\begin{lstlisting}[language=Python]
    os.environ["CUDA_VISIBLE_DEVICES"] = "0"

    HF_CACHE = Path.cwd() / ".hf_cache"
    HF_CACHE.mkdir(parents=True, exist_ok=True)

    os.environ["HF_HOME"] = str(HF_CACHE)
    os.environ["HF_HUB_CACHE"] = str(HF_CACHE / "hub")
\end{lstlisting}
Utilizing a local cache confines model weights to the working directory, preventing redundant downloads across system locations. The transcription model is subsequently instantiated on CUDA using \texttt{float16} precision:
\begin{lstlisting}[language=Python]
    MODEL_NAME = "large-v3"
    LANGUAGE = "it"

    model = WhisperModel(
        MODEL_NAME,
        device="cuda",
        device_index=0,
        compute_type="float16",
        download_root=str(HF_CACHE / "hub"),
    )
\end{lstlisting}
Setting \texttt{language="it"} explicitly conditions decoding on Italian. This parameter bypasses automatic language identification and directly aligns model execution with the linguistic profile of the dataset.
Prior to processing, the pipeline recursively retrieves and filters all \texttt{.mp4} files within the input directory, sorting the resulting collection to guarantee deterministic execution order:
\begin{lstlisting}[language=Python]
    def find_videos(folder: Path) -> list[Path]:
        return sorted(
            path
            for path in folder.rglob("*")
            if path.is_file()
            and path.suffix.lower() in VIDEO_EXTENSIONS
        )
\end{lstlisting}
The core transcription workflow relies on \texttt{model.transcribe()}. By directly parsing the audio embedded within the video container, Faster-Whisper eliminates intermediate extraction steps and yields chronologically ordered transcription segments:
\begin{lstlisting}[language=Python]
    segments, _ = model.transcribe(
        str(video_path),
        language=LANGUAGE,
        beam_size=5,
        vad_filter=True,
        condition_on_previous_text=False,
    )
\end{lstlisting}
Three key hyperparameters govern decoding precision and stability:
\begin{itemize}
    \item \texttt{beam\_size=5}: Activates beam-search decoding, evaluating multiple competing hypotheses during generation to balance accuracy and computational overhead.
    \item \texttt{vad\_filter=True}: Enables Voice Activity Detection to strip environmental noise, music, and unvoiced pauses, preventing non-linguistic signals from reaching the ASR engine.
    \item \texttt{condition\_on\_previous\_text=False}: Disables cross-segment context conditioning, isolating individual speech segments and mitigating error propagation across the audio stream.
\end{itemize}
The resulting segment iterator is aggregated into a continuous transcript by stripping whitespace, discarding empty segments, and concatenating valid text chronologically:
\begin{lstlisting}[language=Python]
    return " ".join(
        segment.text.strip()
        for segment in segments
        if segment.text.strip()
    )
\end{lstlisting}
Omitting downstream linguistic post-processing or semantic filtering preserves an unadulterated representation of raw ASR performance.
The corpus is processed sequentially, using each video filename as a unique primary key:
\begin{lstlisting}[language=Python]
    transcriptions = {}

    for index, video_path in enumerate(videos, start=1):
    try:
        transcription = transcribe_video(video_path)
        transcriptions[video_path.name] = transcription
    except Exception as error:
        transcriptions[video_path.name] = (
            f"ERROR: {error}"
        )
\end{lstlisting}
Robust error handling isolates processing failures on individual files, preventing execution halts and cleanly decoupling unvoiced media from runtime exceptions.
Finally, extracted transcripts are exported to a structured CSV file indexed by \texttt{video\_name} and \texttt{transcription}:
\begin{lstlisting}[language=Python]
    writer = csv.DictWriter(
        csv_file,
        fieldnames=[
        "video_name",
        "transcription"
        ],
        delimiter=";",
    )

    writer.writeheader()

    writer.writerows(
        {
            "video_name": video_name,
            "transcription": transcription,
        }
        for video_name, transcription
        in transcriptions.items()
    )
\end{lstlisting}

%Non Corretto